\chapter{Mathematics}
\label{ch:math}
\section{Foliations vs Fibrations} %gemini
\label{sec:math-fvsf}

While both fibrations and foliations are mathematical concepts used in
differential geometry and topology to decompose complex spaces into simpler,
lower-dimensional pieces, they approach this "slicing" in fundamentally
different ways.
The easiest way to understand the distinction is that a fibration is a strictly
organized, global structure, while a foliation is a purely local slicing rule
that can lead to wild and complex global behaviors.
Here is a breakdown of the differences.

\paragraph{Fibration: The Global Bundle}
A fibration (specifically a fiber bundle) is a way of projecting a large space
down to a smaller space, such that every point in the smaller space corresponds
to an identical "fiber" in the large space.

 You have a total space $ E $, a base space $B$, and a continuous map $p:E
 \to B$. If you pick any point in the base space $B$, its preimage (the
 "fiber" above it) looks exactly like a standard space $F$. Furthermore,
 locally, the total space must look exactly like the product of the base and
 the fiber ( $B \times F$). Every single fiber in the entire space is 
 topologically identical (homeomorphic).

Analogy: Think of a hairbrush. The base space is the plastic handle. The fibers are the bristles. Every bristle is identical, and they are organized perfectly over the base.

Classic Example: The Möbius strip is a fibration (a fiber bundle). The base space 
$B$ is a circle, and the fiber $F$ is a line segment. As you travel around the
circle, the line segments twist, but every single vertical "slice" of the
Möbius strip is an identical line segment.

\paragraph{Foliation: The Local Layers}
A foliation is a way of completely filling an 
$n$-dimensional manifold with a set of disjoint, connected, 
$p$-dimensional submanifolds called leaves.
 There is no global map to a base space. Instead, a foliation is defined by
 local coordinate charts.
If you zoom in anywhere on the manifold, it must look like a stack of parallel
planes (specifically, it locally looks like $\mathbb{R}^p \times \mathbb{R}^{n-p}$, where the 
$\mathbb{R}^p$ slices are the leaves).

 Because the rule is only local, the global behavior of the leaves can be
 chaotic. Leaves do not have to be topologically identical to one another. Some
 leaves might be compact (like a sphere or a torus), while a neighboring leaf
 might be infinitely long and wrap around the compact leaf forever without ever
 closing in on itself.

Analogy: Think of the flaky layers inside a croissant, or the grain lines in a
block of wood. Locally, it just looks like stacked parallel layers. Globally,
those layers can swirl, merge, or spiral in ways that cannot be mapped neatly
to a 2D base.

Classic Example: The Reeb foliation of a 3-dimensional sphere. It slices the 3D
sphere into 2D leaves. One single leaf is a closed torus (a donut shape)
sitting inside the sphere. Every other leaf is an infinite, non-compact flat
plane that spirals around and endlessly approaches the surface of that torus.


\begin{table}[h]
    \centering
    \caption{The Key Differences Summarized}
    \begin{tabular}{p{0.2\textwidth} p{0.4\textwidth} p{0.4\textwidth}}
        \toprule
        Feature	&Fibration	&Foliation\\
        \midrule
        Defining characteristic&	A global projection map to a base space. &	A local decomposition into parallel slices.\\
        Uniformity&	All fibers are topologically identical. &	Leaves can be drastically different from one another globally.\\
        Base Space&	Exists. The fibers are parameterized by a base manifold.&	Does not exist. You cannot generally create a sensible "space of leaves".\\
        Relationship&	Every locally trivial fibration automatically creates a foliation (the fibers serve as the leaves).	&Most foliations do not come from fibrations, because their leaves are not uniform and spiral chaotically.\\
  \bottomrule    
  \end{tabular}  
\end{table}


In short, if a space is a fibration, it is highly organized and every slice is
the same. If a space has a foliation, it merely means it is made of layers that
never intersect, but those layers are free to stretch, spiral, and differ from
one another across the whole space.

\section{FT on manifold}
\label{sec:math-ftm}

There is a fundamental distinction between Fourier analysis in Euclidean space
and harmonic analysis on manifolds. It becomes particularly important when
constructing Wigner functions, Weyl quantization, or phase-space descriptions
of quantum mechanics.

The central idea is that a manifold does not possess a global linear structure,
but each of its tangent spaces does. Consequently, Fourier transformation can
be defined naturally on individual tangent spaces, while the geometry tells us
how these local transforms are related.

There are three principal ways to generalize Fourier analysis:
\begin{enumerate}
   \item Spectral Fourier analysis: replace plane waves by eigenfunctions of a
       geometrically defined operator, usually the Laplace–Beltrami operator.
   \item Tangent-space Fourier analysis: perform ordinary Fourier
       transformation in \(T_xM\), with momenta belonging to \(T_x^*M\). This
        is especially relevant to Wigner functions.
   \item  Group-theoretic Fourier analysis: when the manifold has sufficient
       symmetry, use representations of its symmetry group.

\end{enumerate}
These constructions are related, but they answer different physical questions.

\subsection{What exactly fails on a manifold?}
In Euclidean space,
\[
\widetilde f(p)=\int_{\mathbb R^n}
f(x)e^{-ip\cdot x/\hbar}\,d^nx.
\]
Several structures are hidden in this familiar expression.
First, all points are identified with elements of a single vector space. 
Second, the quantity
\[
p\cdot x
\]
is a canonical pairing of a vector with a covector. 
Third, translations form an Abelian group,
\[
x\mapsto x+a,
\]
whose irreducible unitary representations are the plane waves
\[
\chi_p(a)=e^{ip\cdot a/\hbar}.
\]
Indeed,
\[
\chi_p(a+b)=\chi_p(a)\chi_p(b).
\]
Thus Fourier analysis is fundamentally connected to the representation theory of translations.

On a general manifold \(M\), however:
\begin{enumerate}
   \item There is no canonical difference \(x-y\) between two points.
   \item The cotangent spaces \(T_x^*M\) and \(T_y^*M\) are different vector
       spaces, with no canonical identification.
   \item There is generally no transitive Abelian translation group.
\end{enumerate}
This means that a globally defined Fourier transform with a single momentum
variable \(p\in\mathbb R^n\), behaving exactly as in Euclidean space, is not
naturally available.

But notice something important:
\[
\boxed{
T_xM\text{ is linear, and }T_x^*M
\text{ is its canonical dual.}
}
\]
The pairing
\[
\langle p,v\rangle=p_\mu v^\mu,
\qquad
v\in T_xM,\quad p\in T_x^*M,
\]
is perfectly well defined and coordinate independent.
We therefore have a natural Fourier transform at every point of the manifold,
without requiring a global identification of tangent spaces.

\subsection{ A geometric Fourier transform using the exponential map}
Let \(M\) be an \(n\)-dimensional Riemannian manifold with metric \(g\) and
Levi-Civita connection.
At every \(x\in M\), the exponential map provides a relation between tangent
vectors and points of the manifold:
\[
\exp_x:T_xM\longrightarrow M.
\]
For \(v\in T_xM\), the point
\[
y=\exp_x(v)
\]
is reached by following the geodesic beginning at \(x\) with initial tangent
\(v\), for affine parameter one.
Locally, we can reverse this relation:
\[
v=\log_x(y)=\exp_x^{-1}(y).
\]
This is the geometric replacement for \(y-x\).

We can consequently define a local Fourier transform centered at \(x\):
\[
\boxed{
\mathcal F_x[f](p)
=
\int_{T_xM}
e^{-i\langle p,v\rangle/\hbar}
\chi_x(v)f(\exp_xv)\,d v_x.
}
\]
Here \(p\in T_x^*M\), \(d v_x\) is the volume measure induced by the metric on
\(T_xM\), and \(\chi_x(v)\) is a cutoff restricting the integral to a region
where the exponential map is one-to-one.
This expression has several desirable properties.
It is coordinate independent: \(\langle p,v\rangle\) is a scalar, and both the
exponential map and the measure are geometrically defined.
It reduces locally to ordinary Fourier analysis when \(M=\mathbb R^n\).
Most importantly, no subtraction of manifold points is required. We subtract,
or Fourier-transform, tangent vectors instead.

One must nevertheless distinguish two constructions. The Fourier transform of a
function \(F(x,v)\) on the tangent bundle \(TM\), taken separately in each
vector space \(T_xM\), is intrinsically defined once a fiber measure is chosen.
The transformation above, which starts with a function \(f\) on \(M\),
additionally uses a connection through \(\exp_x\).
Thus the geometry supplies precisely the structure missing from the ordinary
Fourier formula.

\subsection{Geometric Wigner functions}
This is where the tangent-space approach becomes especially useful.
Recall the standard Wigner function on \(\mathbb R^n\):
\[
W_\psi(x,p)
=
\frac{1}{(2\pi\hbar)^n}
\int_{\mathbb R^n}
e^{-ip\cdot v/\hbar}
\psi(x+v/2)\psi^*(x-v/2)\,d^nv.
\]
The problem is not the Fourier transformation itself. It is the appearance of
\[
x+\frac v2,\qquad x-\frac v2.
\]
These Euclidean operations assume an affine structure.

The geometric replacement
Let \(x\in M\), \(v\in T_xM\), and define
\[
x_+=\exp_x(v/2),\qquad
x_-=\exp_x(-v/2).
\]
These points lie symmetrically along a geodesic passing through \(x\).

For a scalar wavefunction, a locally defined geometric Wigner function is
\[
\boxed{
W_\psi(x,p)=
\frac{1}{(2\pi\hbar)^n}
\int_{T_xM}
e^{-i\langle p,v\rangle/\hbar}
\psi(x_+)\psi^*(x_-)
A(x,v)\chi_x(v)\,dv_x.
}
\]
The ingredients are geometric:
\[
x\in M,\qquad
v\in T_xM,\qquad
p\in T_x^*M.
\]
The real, even factor \(A(x,v)\), normalized by \(A(x,0)=1\), allows for
geometric density corrections. The cutoff \(\chi_x\), also even and equal to
one at the origin, confines the construction to a normal neighborhood.

For a normalized scalar wavefunction with \(A(x,0)=1\), integration over
momentum gives the correct local position density:
\[
\int_{T_x^*M}W_\psi(x,p)\,dp_x
=
|\psi(x)|^2.
\]
This follows immediately because
\[
\frac{1}{(2\pi\hbar)^n}
\int_{T_x^*M}
e^{-i\langle p,v\rangle/\hbar}\,dp_x
=
\delta_x(v).
\]
Here the dual momentum measure is normalized consistently with the
tangent-space measure.

However, the position marginal alone does not determine the correct geometric
Wigner transform. Additional requirements, such as invertibility,
self-adjointness of Weyl quantization, and an appropriate Hilbert–Schmidt
identity, distinguish different prescriptions for \(A\).

A particularly useful treatment is the balanced geodesic Weyl calculus
developed by Dereziński, Latosiński and Siemssen. Their construction uses
half-densities and the Van Vleck–Morette determinant, rather than simply
replacing Euclidean displacements with geodesic displacements. 

\paragraph{Why curvature enters?}
The Van Vleck–Morette determinant can be expressed through Synge's world function
\[
\sigma(x,y)=\frac12 d(x,y)^2
\]
in a sufficiently small geodesically convex neighborhood:
\[
\Delta(x,y)=
\frac{
\left|\det[-\nabla_\mu\nabla_{\nu'}\sigma(x,y)]\right|
}{
\sqrt{|g(x)|}\sqrt{|g(y)|}
}.
\]
The primes distinguish differentiation at \(y\) from differentiation at \(x\).
This quantity measures how geodesic separation and volume elements differ from
their Euclidean counterparts. It equals one on the diagonal and in flat
Euclidean space, but generally receives curvature corrections.
In this way, curvature enters geometric Fourier analysis not by changing the
elementary exponential \(e^{-i\langle p,v\rangle/\hbar}\), but through the
relation between tangent-space displacements and actual manifold geometry.

\subsection{Why this construction is not automatically global}
There remains a fundamental difficulty.
The exponential map need not be globally one-to-one. Two distinct tangent
vectors may arrive at the same point, and two manifold points may be joined by
multiple geodesics.

Consider the sphere \(S^2\). Between antipodal points there are infinitely many
minimizing geodesics. Which one should determine the midpoint and the
displacement vector of a Wigner transform?
There is no geometrically distinguished answer.
A local construction works in geodesically convex neighborhoods. A global
construction requires additional choices, restrictions, or a different
formulation.

There are special manifolds for which this difficulty disappears. For example,
on a complete, simply connected Riemannian manifold of nonpositive sectional
curvature, the Cartan–Hadamard theorem guarantees that every exponential map is
a global diffeomorphism. Unique geodesics then make a global midpoint
construction possible.
Even there, the resulting transform is not necessarily equivalent to ordinary
Euclidean Fourier analysis, since curvature remains present.
A key distinction is therefore:
Coordinate independence does not imply the existence of a unique global Fourier
transform.

\subsection{Spectral Fourier analysis: a different solution}
Suppose \(M\) is a compact Riemannian manifold without boundary.
The Laplace–Beltrami operator has an orthonormal eigenbasis:
\[
-\Delta_g\phi_k=\lambda_k\phi_k.
\]
We may define
\[
\widetilde f_k
=
\int_M f(x)\phi_k^*(x)\,d\mu_g(x)
\]
and reconstruct
\[
f(x)=\sum_k\widetilde f_k\phi_k(x),
\]
with convergence in \(L^2(M)\).
This is a genuine global generalization of Fourier series.

For a simple physical example, consider a particle constrained to move on a
sphere of radius \(R\).
The appropriate global basis consists of spherical harmonics:
\[
\phi_{\ell m}=Y_{\ell m}(\theta,\varphi),
\]
and
\[
-\Delta_{S^2}Y_{\ell m}
=
\frac{\ell(\ell+1)}{R^2}Y_{\ell m}.
\]
Consequently, the kinetic-energy eigenvalues are
\[
E_\ell=\frac{\hbar^2}{2mR^2}\ell(\ell+1).
\]
We have a perfectly good global Fourier-like representation without a global
linear configuration space.
But notice the difference: the labels \(\ell,m\) do not form an ordinary
momentum vector, and the degeneracy label \(m\) requires a choice of
angular-momentum axis.

The spectral expansion gives a global decomposition into modes. It does not
directly describe a local momentum covector at a given point.
This distinction is central in phase-space quantum mechanics.

For further discussion, Feichtinger, Führ and Pesenson's Geometric
Space-Frequency Analysis on Manifolds develops the relationship between
spectral frequencies and localization on Riemannian manifolds.

\subsection{Group Fourier transforms}
There is another possibility when \(M\) is itself a Lie group, or when it is a
homogeneous space.
For a compact Lie group \(G\), the Fourier transform is naturally formulated in
terms of irreducible unitary representations \(\pi\):
\[
\widehat f(\pi)
=
\int_G f(g)\pi(g)^\dagger\,dg.
\]
Unlike the Abelian Fourier transform, \(\widehat f(\pi)\) is generally
matrix-valued.
The Peter–Weyl theorem supplies the corresponding completeness and inversion
theory.

The circle \(S^1\) is particularly instructive. Although it does not have the
global linear structure of \(\mathbb R\), it is an Abelian Lie group, and
ordinary Fourier series work globally:
\[
f(\theta)=\sum_{m\in\mathbb Z}f_m e^{im\theta},
\]
\[
f_m=\frac{1}{2\pi}
\int_0^{2\pi}f(\theta)e^{-im\theta}\,d\theta.
\]
The momentum spectrum is discrete:
\[
L_z=m\hbar.
\]
This shows that the absence of global linear structure is not, by itself, fatal
to global Fourier analysis. An appropriate group structure may replace it.

The Wigner function for a quantum rotor must accordingly be adapted to angular
position and angular momentum. Its classical phase space is the cylinder
\(T^*S^1\), not the Euclidean plane. 

\subsection{The deeper point: phase space remains global}
Perhaps the most important observation for classical and quantum mechanics is
that the cotangent bundle
\[
T^*M
\]
exists globally, regardless of whether \(M\) has a global linear structure.
It possesses a canonical one-form
\[
\vartheta=p_i\,dx^i
\]
and canonical symplectic form
\[
\omega=-d\vartheta=dx^i\wedge dp_i.
\]
Although these expressions use local coordinates, the underlying differential
forms are globally defined.
Thus the classical phase space is perfectly legitimate without a global
identification
\[
T^*M\cong M\times\mathbb R^n.
\]

This suggests a particularly useful physical viewpoint.
The momentum of a particle at \(x\) is not an element of some universal
momentum vector space. It is a covector
\[
p\in T_x^*M.
\]
If a semiclassical wavefunction is locally written as
\[
\psi(x)=a(x)e^{iS(x)/\hbar},
\]
then its local momentum is precisely
\[
p=dS.
\]
The phase function \(S\) need not be globally single-valued, and a single
smooth phase may break down at caustics. The cotangent bundle nevertheless
remains well defined.

This also explains an advantage of the Wigner approach: it starts from the
two-point density kernel,
\[
\rho(x,y)=\psi(x)\psi^*(y),
\]
rather than presupposing a globally smooth phase \(S\).
Multiple semiclassical branches and interference can therefore be described
without requiring a single-valued Hamilton–Jacobi action everywhere.

\subsection{An additional geometric ingredient: parallel transport}
For scalar wavefunctions, the product \(\psi(x_+)\psi^*(x_-)\) makes sense
mathematically. But difficulties appear if wavefunctions carry spin or gauge
degrees of freedom.
Suppose, for example, that
\[
\psi(x)\in E_x
\]
where \(E\to M\) is a vector bundle.
Then \(\psi(x_+)\) and \(\psi(x_-)\) belong to different fibers. We cannot
directly construct an endomorphism of a single fiber from them.

The natural remedy is parallel transport.
Let
\[
U(x,y):E_y\longrightarrow E_x
\]
denote parallel transport along the chosen geodesic.
Define transported states
\[
\Psi_\pm(x,v)=U(x,x_\pm)\psi(x_\pm).
\]
Now both vectors belong to \(E_x\), and their outer product
\[
\Psi_+(x,v)\Psi_-(x,v)^\dagger
\]
is an operator on that common fiber.
Fourier transformation of this object over \(T_xM\) gives a matrix-valued,
bundle-covariant Wigner function.

This construction unifies the geometry of the configuration manifold with that
of the internal spin or gauge bundle.
Such tangent-bundle constructions have been developed for quantum fields in
curved spacetime with external Yang–Mills fields; see Fonarev, Wigner Function
and Quantum Kinetic Theory in Curved Space-Time and External Fields, Journal of
Mathematical Physics 35, 2105–2129 (1994). 

\subsection{The most important distinction}
We should not seek a single construction called the Fourier transform on a
manifold.
Instead, there are two different physical objectives.
\paragraph{A global representation of wavefunctions}
Spectral theory, harmonic analysis, and group representations give global mode
decompositions. These are useful for solving wave equations, finding energy
spectra, and constructing propagators.

\paragraph{A local position–momentum description}
Tangent-space Fourier transforms, geodesic Wigner functions, and
pseudodifferential symbols give functions on \(T^*M\). These are useful for
semiclassical dynamics, kinetic theory, and describing local momentum
distributions.

The second objective is particularly close to classical Hamiltonian mechanics.
The central construction can be summarized as
\[
\boxed{
\begin{gathered}
\text{Two-point kernel on }M\times M\\
\big\downarrow\quad \text{geodesic midpoint}\\
\text{Function on }TM\\
\big\downarrow\quad\text{fiberwise Fourier transform}\\
\text{Function on }T^*M.
\end{gathered}
}
\]
This is the natural geometric replacement for the ordinary Wigner transformation.

\section{Musical Isomorphisms}
\label{sec:math-musical}
The term musical isomorphism is a picturesque name from differential geometry
for the correspondence between vectors and covectors induced by a nondegenerate
metric.
The name comes from the two musical symbols:
\begin{itemize}
   \item \(\flat\) (flat), which converts a vector into a covector.
   \item \(\sharp\) (sharp), which converts a covector into a vector.
\end{itemize}
These operations are often called lowering and raising indices.

Their relevance to  Legendre transformations is particularly
important: for a mechanical system with a quadratic kinetic energy, the
Legendre map is precisely a musical isomorphism (possibly multiplied by mass).
Let us examine what this actually means.

\subsection{Vectors and covectors }
Let \(M\) be a smooth manifold and \(q\in M\).
The tangent space \(T_qM\) consists of tangent vectors at \(q\). In
coordinates, a vector is written
\[
v=v^i\frac{\partial}{\partial q^i}.
\]
A covector belongs to the dual space
\[
T_q^*M.
\]
A covector is a linear functional on tangent vectors:
\[
\alpha:T_qM\longrightarrow\mathbb R.
\]
In coordinates,
\[
\alpha=\alpha_i\,dq^i.
\]
Its action on a vector is
\[
\boxed{\alpha(v)=\alpha_iv^i.}
\]
This pairing requires no metric.
The two spaces are dual, but they are not canonically identical.
In particular, there is no natural way to turn an arbitrary vector into a
covector without supplying additional structure.

\subsection{ A metric provides precisely the missing structure}
A metric tensor \(g\) is a nondegenerate, symmetric bilinear form
\[
g:T_qM\times T_qM\longrightarrow\mathbb R.
\]
For a Riemannian metric, \(g\) is positive definite. In relativistic geometry,
\(g\) is generally indefinite.
Given a vector \(v\), we may construct a covector by holding the first argument
of the metric fixed:
\[
\boxed{v^\flat=g(v,\cdot).}
\]
The notation means that
\[
v^\flat(w)=g(v,w)
\]
for every vector \(w\).
In coordinates,
\[
\boxed{(v^\flat)_i=g_{ij}v^j.}
\]
This is called the flat musical isomorphism:
\[
\flat:T_qM\longrightarrow T_q^*M.
\]
Because the metric is nondegenerate, this map is invertible.
Its inverse is the sharp musical isomorphism:
\[
\sharp:T_q^*M\longrightarrow T_qM.
\]
For a covector \(\alpha\),
\[
\boxed{(\alpha^\sharp)^i=g^{ij}\alpha_j,}
\]
where
\[
g^{ik}g_{kj}=\delta^i_j.
\]
Consequently,
\[
\boxed{
(v^\flat)^\sharp=v,
\qquad
(\alpha^\sharp)^\flat=\alpha.
}
\]
The two maps are linear isomorphisms between the tangent and cotangent spaces.

\subsection{Geometric example: Euclidean space}
In Cartesian coordinates on ordinary Euclidean space,
\[
g_{ij}=\delta_{ij}.
\]
Therefore,
\[
v_i=\delta_{ij}v^j=v^i.
\]
So the components of a vector and its associated covector have the same
numerical values.
For example,
\[
v=(2,3)
\]
becomes the covector
\[
v^\flat=2\,dx+3\,dy.
\]
The distinction can easily go unnoticed because the components coincide.
But the geometric meanings remain different.
The vector \(v\) describes a direction and magnitude of displacement.
The covector \(v^\flat\) assigns to any displacement \(w\) the scalar
\[
v^\flat(w)=v\cdot w.
\]
The metric has converted a geometric displacement into a linear functional on
displacements.
What happens in curvilinear coordinates?
Consider the Euclidean plane in polar coordinates:
\[
ds^2=dr^2+r^2d\varphi^2.
\]
The metric components are
\[
g_{ij}=
\begin{pmatrix}
1&0\\
0&r^2
\end{pmatrix}.
\]
For a vector
\[
v=v^r\frac{\partial}{\partial r}
+v^\varphi\frac{\partial}{\partial\varphi},
\]
the associated covector is
\[
\boxed{
v^\flat=v^r\,dr+r^2v^\varphi\,d\varphi.
}
\]
Thus
\[
v_r=v^r,\qquad v_\varphi=r^2v^\varphi.
\]
The vector and covector no longer have the same coordinate components, even
though the underlying geometry is still Euclidean.
This illustrates why the distinction is not merely formal.

\subsection{Connection with the Legendre transformation}
Now consider a mechanical system with configuration manifold \(Q\) and
Lagrangian
\[
L(q,v)=\frac12 m\,g_{ij}(q)v^iv^j-V(q).
\]
The canonical momentum is
\[
p_i=\frac{\partial L}{\partial v^i}.
\]
Differentiating gives
\[
\boxed{p_i=mg_{ij}(q)v^j.}
\]
But the musical isomorphism gives
\[
(v^\flat)_i=g_{ij}v^j.
\]
Therefore,
\[
\boxed{p=mv^\flat.}
\]
The Legendre map is exactly the flat musical isomorphism multiplied by mass.
Its inverse is
\[
\boxed{v=\frac1m p^\sharp.}
\]
The Hamiltonian is then
\[
H(q,p)=\frac{1}{2m}g^{ij}(q)p_ip_j+V(q).
\]
This is the geometric meaning of the familiar Newtonian expression
\[
p=mv.
\]
More precisely, the correct geometric statement is
\[
\boxed{p=m\,g(v,\cdot).}
\]
The equation \(p=mv\) implicitly identifies vectors and covectors using the
Euclidean metric.

Example: a particle in polar coordinates
The Lagrangian is
\[
L=\frac m2(\dot r^2+r^2\dot\varphi^2)-V(r,\varphi).
\]
The canonical momenta are
\[
p_r=m\dot r,
\qquad
p_\varphi=mr^2\dot\varphi.
\]
The inverse relations are
\[
\dot r=\frac{p_r}{m},
\qquad
\dot\varphi=\frac{p_\varphi}{mr^2}.
\]
Hence
\[
\boxed{
H=\frac{p_r^2}{2m}
+\frac{p_\varphi^2}{2mr^2}+V.
}
\]
The factor \(r^2\) arises directly from the metric and its musical
isomorphisms.
Notice that \(p_\varphi\) is angular momentum, not linear momentum. This is
exactly what we expect from a covector conjugate to an angular coordinate.

An important distinction must be made.
The musical isomorphism is linear on every tangent space:
\[
(v+w)^\flat=v^\flat+w^\flat.
\]
The Legendre map of a general Lagrangian,
\[
\mathbb FL(q,v)
=\left(q,\frac{\partial L}{\partial v}\right),
\]
need not be linear.

For example, consider the relativistic Lagrangian of a free particle in a fixed
inertial frame,
\[
L(v)=-mc^2\sqrt{1-\frac{v^2}{c^2}}.
\]
Then
\[
p_i=\frac{mv_i}{\sqrt{1-v^2/c^2}},
\]
where the spatial indices are lowered using the positive Euclidean spatial
metric.
The factor
\[
\gamma(v)=\frac1{\sqrt{1-v^2/c^2}}
\]
depends on velocity, so the map \(v\mapsto p\) is nonlinear.

Thus it is not a musical isomorphism associated with a fixed metric.
The distinction is especially relevant in Finsler geometry, where the kinetic
function need not be quadratic in velocity.

\subsection{ Musical isomorphisms in relativity}
In special relativity with signature
\[
\eta_{\mu\nu}=\operatorname{diag}(+,-,-,-),
\]
consider a four-vector
\[
v^\mu=(v^0,v^1,v^2,v^3).
\]
Its metric dual is
\[
v_\mu=\eta_{\mu\nu}v^\nu,
\]
so
\[
\boxed{
v_\mu=(v^0,-v^1,-v^2,-v^3).
}
\]
The flat operation changes the signs of the spatial components.
For a four-momentum,
\[
p^\mu=(E/c,\mathbf p),
\]
the corresponding covector is
\[
p_\mu=(E/c,-\mathbf p).
\]
This matters whenever we distinguish tangent vectors from momentum covectors on
spacetime.
For example, the free relativistic particle satisfies
\[
p_\mu p^\mu=m^2c^2.
\]
The metric is necessary to form this scalar from the momentum covector.
The canonical symplectic form on spacetime phase space, however,
\[
\omega=dx^\mu\wedge dp_\mu,
\]
does not require the spacetime metric.
That distinction connects directly to our earlier discussion: the metric
identifies vectors with covectors, while the symplectic form determines the
canonical pairing and Hamiltonian dynamics.

\subsection{ Symplectic geometry also provides a musical isomorphism}
The terminology is not restricted to metrics.
Any nondegenerate bilinear form can establish an isomorphism between a vector
space and its dual.
For a symplectic form \(\omega\), define
\[
\boxed{
X^{\flat_\omega}=\iota_X\omega=\omega(X,\cdot).
}
\]
Since \(\omega\) is nondegenerate, this map is invertible:
\[
\sharp_\omega:T^*M\longrightarrow TM.
\]
Hamilton's equations can then be written
\[
\boxed{
X_H=(dH)^{\sharp_\omega}.
}
\]
The metric and symplectic musical maps differ fundamentally.
For a metric,
\[
g(X,X)
\]
measures the squared length (or squared norm, in an indefinite metric) of \(X\).
For a symplectic form,
\[
\omega(X,X)=0
\]
for every \(X\).
Thus their musical isomorphisms have very different physical consequences.
A metric turns \(dH\) into the gradient of energy. A symplectic form turns
\(dH\) into the vector field generating conservative Hamiltonian evolution.

\subsection{Summary}
\begin{table}[h]
    \caption{The essential structures}
    \begin{tabular}{lcr}
        \toprule
        Structure &	Map	&Interpretation\\
        \midrule
        Metric \(g\)&	\(v\mapsto g(v,\cdot)\)	&Vector to metric-dual covector\\
        Inverse metric &	\(\alpha\mapsto\alpha^\sharp\)	&Covector to vector\\
        Quadratic kinetic energy &	\(v\mapsto mg(v,\cdot)\) &	Velocity to canonical momentum\\
        General Lagrangian	&\(v\mapsto\partial L/\partial v\)	&General, possibly nonlinear Legendre map\\
        Symplectic form \(\omega\)&	\(X\mapsto\omega(X,\cdot)\)	&Vector field to covector\\
        Hamiltonian dynamics&	\(dH\mapsto X_H\)&	Energy differential to evolution vector field\\
\bottomrule
    \end{tabular}
\end{table}
The main point is that a vector and a covector are intrinsically different
objects. A musical isomorphism identifies them using a chosen geometric
structure; it does not make them intrinsically identical.
In ordinary Cartesian mechanics, this distinction is almost invisible. On
curved configuration manifolds, in curvilinear coordinates, in relativity, and
in geometric Hamiltonian mechanics, it becomes essential.

One last interesting observation: a metric musical isomorphism is a linear map
\(TQ\to T^*Q\), whereas a regular Legendre map can be nonlinear. Yet the latter
still pulls the canonical symplectic structure on \(T^*Q\) back to \(TQ\). This
is precisely how the symplectic geometry of a Lagrangian system arises without
requiring its Legendre map to be a musical isomorphism.
